% ---------------------------------------------------------------------------
%  02-background.tex -- the minimum zone theory, stated so both readers can use it
% ---------------------------------------------------------------------------
\section{Background: zones, DBMs, and the two ways to stay finite}
\label{sec:background}

This section fixes notation and states the three standard facts the engine
depends on.  A reader who knows \uppaal{} internals can skim to
\cref{sec:bg-finite}; a reader who does not should read the boxes.

\subsection{Zones and DBMs}

Let $\Clocks = \{x_1,\dots,x_n\}$ be the clocks and let $x_0 \equiv 0$ be a
reference clock.  A \emph{zone} is a set of clock valuations definable by a
conjunction of constraints $x_i - x_j \prec c$ with ${\prec} \in \{<,\le\}$.
It is represented by a \emph{difference bound matrix} (DBM): a
$\dim \times \dim$ array $D$, $\dim = n+1$, where $D[i][j]$ encodes the
tightest known bound on $x_i - x_j$.  \tool{} packs the bound and its
strictness into one integer, $\mathit{raw} = (c \ll 1)\,|\,\mathit{weak}$, so a
comparison of two bounds is one integer comparison
(\file{HDBM/include/ptctlmc/dbm\_ops.hpp}).

A DBM is \emph{canonical} when every entry is the weight of a shortest path in
the induced constraint graph; canonicalisation is Floyd--Warshall,
$O(\dim^3)$.  Canonical form matters because it makes syntax match semantics:
for canonical $D, D'$,
\[
  [D] \subseteq [D'] \quad\iff\quad D[i][j] \le D'[i][j] \text{ for all } i,j ,
\]
i.e.\ inclusion of the represented sets is entrywise comparison of the
matrices.

\begin{intuition}
A zone is a shape in ``stopwatch space'' cut out by constraints of the form
``$x$ minus $y$ is at most $3$''.  The matrix is just the table of all such
bounds.  Canonicalising means tightening every entry to the best value
implied by the others, so that comparing two shapes becomes comparing two
tables cell by cell.
\end{intuition}

The operations the search needs are all standard and all in
\file{dbm\_ops.hpp}: conjoin a guard (\code{dbm\_constrain}, incremental, $O(\dim^2)$),
reset a clock (\code{dbm\_reset}, $O(\dim)$), let time pass (\code{dbm\_delay},
drop upper bounds), emptiness (check the diagonal), and inclusion
(\code{dbm\_includes}).

\subsection{The two ways to keep the zone graph finite}
\label{sec:bg-finite}

The exact zone graph of a timed automaton can be infinite, so every tool needs
a finiteness mechanism.  There are two families, and it is important for this
document that they are \emph{not} interchangeable.

\paragraph{Widening (extrapolation).}
Replace each zone $Z$ by a coarser zone $\extraLU(Z) \supseteq Z$ chosen so
that only finitely many images exist.  $\extraLU$ is \uppaal{}'s default: for
each clock $x$ it uses $L_x$, the largest constant $x$ is compared against from
below, and $U_x$, the largest from above, and relaxes any bound that is already
outside those windows.  Widening \emph{maps many zones onto one}.  It is sound
for reachability of discrete states, and for clock constraints whose constants
are within $L,U$~\cite{BBLP06}.  It is implemented in
\code{dbm\_extrapolate\_lu\_plus}.

\paragraph{Pruning (abstract subsumption).}
Keep zones exact, but discard a new zone when a stored one already
\emph{simulates} it.  The relevant simulation is the $LU$-preorder
$\preceq_{LU}$ on valuations, and its induced closure
\[
  \aLU(W) \;=\; \{\, v \;:\; \exists w \in W.\; v \preceq_{LU} w \,\}.
\]
The test used is
\[
  Z \subs Z' \quad:\iff\quad [Z] \subseteq \aLU([Z']),
\]
decided syntactically in $O(\dim^2)$ by \code{dbm\_is\_alu\_le}, a port of
\tchecker{}'s \code{is\_alu\_le}, itself the test of
Herbreteau--Srivathsan--Walukiewicz~\cite{HSW12}.  Pruning \emph{never
enlarges a zone}; it only decides that exploring one is pointless.

\begin{intuition}
Two different ways to make an infinite search finite.  \emph{Widening}: round
every shape outward to one of finitely many ``standard'' shapes -- fast, but
you lose precision and can end up exploring configurations the system cannot
actually reach (still giving correct \emph{answers}, just more work and bigger
state counts).  \emph{Pruning}: keep shapes exact, but throw away a new shape
if an old one already dominates it.  \uppaal{} widens by default.  \tool{}
prunes by default and only widens when it detects that pruning is not paying.
\end{intuition}

The two standard facts we rely on, and do \emph{not} reprove:

\begin{proposition}[standard, \cite{HSW12}]\label{prop:alu-finite}
$\preceq_{LU}$ is a preorder on valuations, $\aLU$ is the induced downward
closure, $\subs$ is a simulation for the symbolic successor operator, and
$\subs$-antichains of zones are finite.  Consequently a search that stores
exact zones and prunes with $\subs$ terminates.
\end{proposition}

\begin{proposition}[standard, \cite{BBLP06}]\label{prop:extra-sound}
$\extraLU$ has finite image, $Z \subseteq \extraLU(Z) \subseteq \aLU(Z)$, and
$\extraLU$ preserves satisfaction of every clock constraint whose constants are
bounded by $L$ and $U$.
\end{proposition}

\cref{prop:alu-finite}'s ``$\aLU$ is a downward closure of a preorder'' is
exactly what \cref{sec:theory} will use, so it is worth spelling out the
consequence we need.

\begin{lemma}[$\subs$ is a preorder]\label{lem:preorder}
For a fixed $(L,U)$, $\subs$ is reflexive and transitive.
\end{lemma}

\begin{proof}
Downward closure under a preorder is a closure operator: it is \emph{extensive}
($W \subseteq \aLU(W)$, by reflexivity of $\preceq_{LU}$), \emph{monotone}
($W \subseteq W' \Rightarrow \aLU(W) \subseteq \aLU(W')$, immediate), and
\emph{idempotent} ($\aLU(\aLU(W)) = \aLU(W)$, by transitivity of
$\preceq_{LU}$).  Reflexivity of $\subs$ is extensivity.  For transitivity,
suppose $Z \subs Z'$ and $Z' \subs Z''$.  Then
$[Z] \subseteq \aLU([Z'])$ and $[Z'] \subseteq \aLU([Z''])$, so by monotonicity
and idempotence
$\aLU([Z']) \subseteq \aLU(\aLU([Z''])) = \aLU([Z''])$, whence
$[Z] \subseteq \aLU([Z''])$.
\end{proof}

\begin{remark}
$(L,U)$ is fixed \emph{per discrete state}: the bounds are computed from the
location vector, and every zone in one antichain belongs to the same discrete
state by construction (the antichain is indexed by $\dk$, which encodes the
location vector and the integer variables losslessly).  So \cref{lem:preorder}
applies within an antichain, which is where every use of transitivity in
\cref{sec:theory} occurs.  This is not a technicality one may skip: if the
bounds varied inside a chain, transitivity would fail and the concurrency
argument with it.
\end{remark}

Finally, one cheap implication we will fire constantly:

\begin{lemma}[inclusion implies subsumption]\label{lem:incl-implies-alu}
If $Z$ and $Z'$ are canonical and $[Z] \subseteq [Z']$, then $Z \subs Z'$.
\end{lemma}

\begin{proof}
Semantically this is extensivity of $\aLU$.  Syntactically -- which is what the
code relies on -- \code{dbm\_is\_alu\_le}$(Z,Z')$ returns false only if it finds
a witness pair $(x,y)$ with $Z'[y][x] < Z[y][x]$.  Canonical inclusion gives
$Z[i][j] \le Z'[i][j]$ for all $i,j$, so no such witness exists.
\end{proof}

\begin{remark}[why \cref{lem:incl-implies-alu} matters here]
It licenses the engine to run a cheap plain-inclusion test \emph{first} and, on
a hit, skip the $\aLU$ test entirely.  Since the cheap test can be evaluated
directly on the compressed zone (\cref{lem:mingraph}), this is what removes
decompression from the hot path -- and therefore from inside the lock.
\end{remark}
